
\documentclass[10pt,a4paper]{article}
\usepackage{fontspec}
\usepackage[russian]{babel}
\setmainfont{Times New Roman} 
\usepackage[left=1.8cm,right=1.8cm,top=1.8cm,bottom=1.8cm]{geometry}
\usepackage{hyperref}
\usepackage{xcolor}
\usepackage{enumitem}


\definecolor{primary}{RGB}{128, 0, 32}  
\definecolor{accent}{RGB}{236, 64, 122}     
\definecolor{subtext}{RGB}{180, 80, 110}      


\hypersetup{
    colorlinks=true,
    linkcolor=accent,
    filecolor=accent,      
    urlcolor=accent,
}


\setlist[itemize]{leftmargin=*,itemsep=2pt,topsep=2pt}
\pagestyle{empty}
\usepackage{tikz}

\newcommand{\linedsection}[1]{%
  \vspace{10pt}%
  \noindent
  {\bfseries\Large\color{primary} #1}%
  \par\vspace{3pt}%
  {{\color{primary}\hrule height 1pt}}
  \par\vspace{6pt}%
  \noindent
}



\begin{document}

% --- ШАПКА И КОНТАКТЫ ---
\begin{center}
    {\Huge \bfseries \color{primary} Кира Кошелева} \\[6pt]
    {\large Студентка 2 курса · Специальность 14.05.04 «Электроника и автоматика физических установок»} \\[4pt]
    \small 
    Email: \href{mailto:kosheleva.kir@yandex.ru}{kosheleva.kir@yandex.ru} \quad$\cdot$\quad 
    Telegram: \href{https://t.me/koshelz}{@koshelz} \quad$\cdot$\quad 
    GitHub: \href{https://github.com/koshelz}{koshelz} \quad$\cdot$\quad 
    GitLab: \href{https://gitlab.com/koshelz}{koshelz} 
\end{center}

\vspace{-6pt}
{\color{primary}\hrule height 1pt}
\vspace{6pt}

% --- ОБРАЗОВАНИЕ ---
\linedsection{ОБРАЗОВАНИЕ}
\vspace{2pt}
\textbf{\large НИЯУ МИФИ} \hfill \textcolor{subtext}{2025 -- 2031 (в процессе)} \\
\textit{Направление: Технологии наноэлектронных приборов и архитектура вычислительных систем} \\
\small {Средний балл:} 5.0 / 5.0

% --- ТЕХНИЧЕСКИЕ НАВЫКИ ---
\linedsection{ТЕХНИЧЕСКИЕ НАВЫКИ}
\begin{itemize}
    \item \textbf{Языки программирования:} C/C++, Python (реализация алгоритмов, обработка данных).
    \item \textbf{Инструменты и ПО:} LaTeX, VS Code.
    \item \textbf{Иностранные языки:} Русский (родной), Английский (чтение технической документации).
\end{itemize}

% --- ПРОЕКТНАЯ ДЕЯТЕЛЬНОСТЬ ---
\linedsection{ПРОЕКТНАЯ ДЕЯТЕЛЬНОСТЬ}
\vspace{2pt}
\textbf{\large Станок автоматической для намотки катушек} \hfill \textcolor{subtext}{2026} \\
\textit{Роль: Инженер-программист встроенных систем} \quad $\cdot$ \quad \href{https://gitlab.com/sstund/winding-machine}{\textbf{GitLab}}
\begin{itemize}
    \item Проектирование и отладка ПО: объединение аппаратных узлов (приводы, интерфейс) и веб-интерфейса в единую систему.
    \item Разработка веб-сервиса удалённого управления для беспроводного ввода параметров намотки через локальную сеть.
    \item Результат: Призовое место (\href{https://gitlab.com/sstund/winding-machine/-/blob/main/2026-05-15_001.jpg}{\color{accent}Диплом призера}) на научной конференции.
\end{itemize}


% --- НАУЧНАЯ ДЕЯТЕЛЬНОСТЬ И ПУБЛИКАЦИИ ---
\linedsection{НАУЧНАЯ ДЕЯТЕЛЬНОСТЬ И ПУБЛИКАЦИИ}
\vspace{2pt}
{\large \textbf{Наноэлектронные датчики для мониторинга элементов железобетонных конструкций}} \hfill \textcolor{subtext}{2026} \\
\small{Научно-аналитический обзор · Участие с докладом на конференциях:}
\begin{itemize}
    \item \href{https://github.com/koshelz/my-CV/blob/main/%D0%A1%D0%B5%D1%80%D1%82%D0%B8%D1%84%D0%B8%D0%BA%D0%B0%D1%82%20%D0%92%D0%BE%D0%93%D0%A3.pdf}{\color{accent}Международная научная конференция «Молодые исследователи– регионам»}
    \item \href{https://github.com/koshelz/my-CV/blob/main/%D0%A1%D0%B5%D1%80%D1%82%D0%B8%D1%84%D0%B8%D0%BA%D0%B0%D1%82%20%D1%81%20%D0%92%D0%9D%D0%A1%20%D0%A1%D0%9D%D0%9E.pdf}{\color{accent} Всероссийская весенняя научная сессия СНО НИЯУ МИФИ}
\end{itemize}



\end{document}